\documentclass[aspectratio=169,11pt]{beamer}
\usetheme[academique]{guyeux}
\usepackage[french]{babel}
\usepackage{graphicx}
\graphicspath{{./}}
\title{Le pathog\`ene comme archive sociale}
\subtitle{\textit{M.~africanum} et la carte ethnolinguistique ouest-africaine}
\author{Christophe Guyeux}
\institute{FEMTO-ST, UMR 6174 CNRS \textperiodcentered\ Universit\'e Marie et Louis Pasteur}
\date{9 septembre 2026}
\begin{document}
\begin{frame}[plain]\titlepage\end{frame}

\section{Le r\'esultat central}

\begin{frame}{Les lign\'ees \'epousent la partition ethnolinguistique}
\slidefigure{demo_fig.pdf}{Le chromosome Y s\'epare le nord du sud ; l'ADN mitochondrial est plat partout.}
\end{frame}

\begin{frame}{Le chromosome Y raconte la m\^eme histoire}
\bigchiffre{$\times$15}{Le chromosome Y est quinze fois plus structur\'e que l'ADN mitochondrial.}
\end{frame}

\begin{sliderupture}
Le pathog\`ene est une patriline sociale.
\end{sliderupture}

\begin{frame}{Deux \'echelles emboît\'ees, pas une contradiction}
\slidecompare{\'Echelle grossi\`ere}{%
  Pourquoi \textit{M.~africanum} est en Afrique de l'Ouest.\par\vskip1.5mm
  \textbf{Ascendance g\'en\'etique} ouest-africaine. Fronti\`ere ancienne et profonde.}{\'Echelle fine}{%
  Quelle sous-lign\'ee, L5 ou L6, on porte.\par\vskip1.5mm
  \textbf{Barri\`ere ethnolinguistique}. R\'ecente et sociale.}
\kb{On \'etend Asante-Poku et Gagneux, on ne les contredit pas.}
\end{frame}

\begin{frame}{Ce que la mesure dit exactement}
\centering
\begin{tabular}{lrr}\toprule
Lign\'ee & Souches & $F_{ST}$ inter-ethnies \\ \midrule
L5 & 812 & 0{,}006 \\
L6 & 1183 & 0{,}005 \\ \bottomrule
\end{tabular}
\cle{Les peuples ouest-africains sont g\'en\'etiquement quasi identiques : la structure fine ne peut pas \^etre porte\'e par la g\'en\'etique profonde.}
\notion{$F_{ST}$}{Part de la diversit\'e g\'en\'etique attribuable \`a la diff\'erenciation entre groupes plut\^ot qu'\`a la variation interne. Proche de 0 : groupes indistinguables.}
\end{frame}

\begin{frame}{Une liste, quand rien de mieux n'existe}
\begin{itemize}
  \item Association forte lign\'ee-ethnie
  \item Isolement par la distance contr\^ol\'e
  \item Cophylog\'enie concordante
  \item T\'emoin n\'egatif : \textit{M.~ulcerans}
\end{itemize}
\kb{Batterie de tests concordants : tous positifs.}
\end{frame}

\appendix
\begin{frame}{Annexe : d\'etail des tests}
Slide back-pocket, hors du compte des slides du corps.
\end{frame}
\end{document}
